\section{Moving from PhyloNet to \phynetpy}
\label{sec:comparison}

This section is addressed to researchers with working PhyloNet analyses. The
useful question is not which tool is better but what specifically changes, and
what does not carry over.

\begin{figure}[t]
  \centering
  \includegraphics[width=\linewidth]{transition.pdf}
  \caption{The same body of methods under two organisations. On the left each
  analysis is a named command that handles its own input, search and scoring.
  On the right the analyses are cells of a product of three axes over shared
  machinery, so the empty cells are visible and distinguishable: $-$ marks a
  combination that is well defined but unbuilt, $\times$ one that is not
  defined for that data type. The figure makes no claim that one organisation
  produces better inferences.}
  \label{fig:transition}
\end{figure}

\begin{table}[t]
  \centering
  \scriptsize
  % Comparison of PhyloNet and PhyNetPy.
%
% Hand-written, unlike the numeric tables, because the rows are qualitative.
% Every PhyNetPy entry is verified against the 0.6.0 source; see INVENTORY.md.
% PhyloNet entries describe version 3.8.2 as distributed and are drawn from its
% published description and its command set.
%
% Rows marked [TODO] could not be verified from the material available and must
% be checked against PhyloNet's documentation before submission.

\begin{tabular}{@{}p{38mm}p{52mm}p{52mm}@{}}
\toprule
& \textbf{PhyloNet 3.8.2} & \textbf{\textsc{PhyNetPy} 1.0.0} \\
\midrule

Implementation language
& Java
& Python, with four Cython extension modules \\

Invocation
& A NEXUS file containing a \texttt{PHYLONET} block, run through
  \texttt{java -jar}
& Python function calls; no driver file \\

Programmatic API
& Java classes are reachable by writing Java against the jar
& The documented surface is the package itself: 121 exported names, two
  inference verbs plus \texttt{simulate} \\

Analysis organisation
& One named command per analysis
& A method is a (data, model, criterion) cell; dispatch goes through a
  registry that also reports validity \\

Adding a method
& Implement a new command
& Register one cell with \texttt{@register}; search drivers, moves,
  parameter optimiser and I/O are inherited \\

Network representation
& Internal to the tool
& \texttt{Network}: a reusable rooted DAG, independent of any analysis, with
  no cap on reticulation in-degree \\

Distinguishing invalid from unbuilt requests
& An unsupported combination surfaces as a command-level error
& \texttt{TypeError}, \texttt{ValueError} and \texttt{NotImplementedError} are
  raised for three distinct causes \\

\midrule
\multicolumn{3}{@{}l}{\emph{Analyses}} \\

Maximum likelihood from gene trees
& \texttt{InferNetwork\_ML}
& Available \\

Maximum pseudo-likelihood from gene trees
& \texttt{InferNetwork\_MPL}
& Available \\

Bayesian from gene trees
& \texttt{MCMC\_GT}
& Available \\

Bayesian from sequences
& \texttt{MCMC\_SEQ}
& Available \\

Marker likelihood
& \texttt{MLE\_BiMarkers}
& Scoring available; the search path raises \\

Bayesian from markers
& \texttt{MCMC\_BiMarkers}
& Available \\

Maximum parsimony, diploid hybridization
& \texttt{InferNetwork\_MP}
& Not in 0.6.0; scheduled for 1.0.0 \\

Maximum parsimony under polyploidy
& \texttt{InferNetwork\_MP\_Allopp}
& Available, with bootstrapping \\

\midrule
\multicolumn{3}{@{}l}{\emph{Supporting capability}} \\

Simulation
& Available as separate commands
& \texttt{simulate} inverts the same axes and returns a data object that feeds
  \texttt{infer} directly; gene trees, sequences and markers under
  \texttt{MSC} only \\

Network comparison
& Distance commands are provided
& Nine dissimilarity measures plus a reticulation-specific comparison module
  (tripartitions, precision and recall) \\

Chain diagnostics
& Output is written for external tools
& ESS, autocorrelation time, HPD intervals, Geweke and Gelman--Rubin
  in-process; Tracer and NEXUS output for external tools \\

Visualisation
& \textbf{[TODO: verify what PhyloNet 3.8.2 provides]}
& ASCII rendering only \\

Across-site rate heterogeneity
& \textbf{[TODO: verify]}
& Not implemented \\

\midrule
\multicolumn{3}{@{}l}{\emph{Distribution and testing}} \\

Installation
& Download a jar; requires a JVM
& \texttt{pip install phynetpy}; requires a C compiler when building from
  source \\

Automated testing
& \textbf{[TODO: verify what PhyloNet publishes]}
& 1{,}130 tests, 1{,}125 passing in the default run \\

Continuous integration
& \textbf{[TODO: verify]}
& GitHub Actions across Python 3.9--3.14, plus Windows and packaging jobs \\

Numerical cross-validation
& ---
& 14 cases checked against PhyloNet itself to a tolerance of $10^{-5}$ \\

License
& \textbf{[TODO: verify PhyloNet's license]}
& MIT \\

\bottomrule
\end{tabular}

  \caption{PhyloNet 3.8.2 and \phynetpy\ 1.0.0. \phynetpy\ entries are verified
  against the source; entries marked \textbf{[TODO]} could not be established
  from the material available and must be checked against PhyloNet's
  documentation before submission rather than guessed.}
  \label{tab:comparison}
\end{table}

\subsection{What changes}

\paragraph{The driver file disappears.}
A PhyloNet analysis is a NEXUS file with a \code{PHYLONET} block, executed
through \code{java -jar}. Parameters are positional and flag-based inside that
block, and a sweep over conditions means generating files. In \phynetpy\ the
analysis is Python: a loop over $\theta$ values or reticulation budgets is a
loop, and intermediate results stay as objects.

\paragraph{Choosing a method becomes choosing arguments.}
Switching from maximum likelihood to pseudo-likelihood changes one argument
rather than one command name, and the return type does not change with it. In
PhyloNet the corresponding change is to a different command with its own syntax
and output format.

\paragraph{Results are objects, not text to parse.}
PhyloNet writes inferred networks and scores to standard output, and a pipeline
around it is usually a parser. \code{infer} returns an \code{InferenceResult}
whose \code{.best} is a \code{Network} that the distance functions, the
decomposition routines and \code{score} all accept directly. This is the
practical difference most PhyloNet users will notice first.

\paragraph{Unsupported requests say which kind of unsupported they are.}
The three-way distinction between \code{TypeError}, \code{ValueError} and
\code{NotImplementedError} tells a user whether to change the data, change the
criterion, or wait.

\subsection{What does not carry over}

\paragraph{Input files are not automatically compatible.}
\phynetpy\ reads extended Newick, NEXUS, FASTA and VCF, so the \emph{data} in a
PhyloNet NEXUS file is usually readable. The \code{PHYLONET} command block is
not: there is no translator from a PhyloNet block to a \phynetpy\ call, and the
migration is a rewrite of the driver, not a conversion.

\paragraph{Results will not be identical.}
The two implementations agree on the likelihood of a fixed network to $10^{-5}$
(Section~\ref{sec:validation}), but searches are stochastic and their proposal
kernels differ, so inferred networks can differ --- and in
Section~\ref{sec:headtohead} they did on one reticulate condition. The
\code{phylonet} search preset exists to narrow this gap when the point is a
comparison rather than an analysis, by reproducing PhyloNet's
optimise-every-parameter-per-topology behaviour.

\paragraph{Some analyses are not available.}
Maximum parsimony from gene trees under diploid hybridization is absent, and
marker maximum likelihood can score but not search. Anyone whose work depends on
those should continue to use PhyloNet.

\paragraph{Installation requirements differ.}
PhyloNet needs a JVM and a jar. \phynetpy\ installs with \code{pip}, but the
compiled graph core is required, so installing from a source distribution needs
a C compiler. Where a wheel is published for the target platform this does not
arise.

\subsection{A worked correspondence}

A PhyloNet pseudo-likelihood analysis reads, inside a NEXUS file:

\begin{lstlisting}[language={}]
BEGIN PHYLONET;
InferNetwork_MPL (gt1,gt2,...,gt100) 1 -pl 4;
END;
\end{lstlisting}

The corresponding \phynetpy\ analysis:

\begin{lstlisting}
from phynetpy.data import GeneTrees
from phynetpy.criteria import PseudoLikelihood
from phynetpy.infer import infer

gts = GeneTrees.from_file("genes.nex", {"A": ["A1", "A2"], "B": ["B1"]})
result = infer(gts, criterion=PseudoLikelihood(), max_reticulations=1)
\end{lstlisting}

The species-to-allele mapping moves from a \code{-a} flag into the data object,
where it travels with the data it describes instead of being restated per
command. The reticulation budget is a keyword rather than a positional integer.
The gene trees are named once, by file, rather than enumerated in the call.

These two are not equivalent in every respect, and the differences are the
point of this section rather than a caveat to it: the searches differ, the
defaults differ, and \code{-pl 4} requests four-way parallelism from PhyloNet
whereas the \phynetpy\ call above is single-threaded. A migration should
re-establish results on data with a known answer rather than assume the two
calls are interchangeable.
